\section{Neural Network Customization and NNI Code Generation}
\label{sec:customization}

Once a base model has been discovered and profiled, users may require architectural modifications to meet specific project constraints or improve performance for their target task. Workflow 5 implements an autonomous customization pipeline that combines \textit{RAG-improved best practice retrieval}, \textit{natural language intent parsing}, and \textit{LLM-driven code generation} to produce executable NNI hyperparameter optimization experiments without manual scripting.

\subsection{Model Architecture Inspection}

The customization workflow begins with detailed introspection of the downloaded model file via the \texttt{inspect\_model\_architecture()} node. This analysis extracts several key features, including the layer-by-layer architectural details such as layer types (Conv2D, Dense, Dropout), their parameters, and output shapes. It also identifies the expected input and output tensor shapes along with their data types, calculates the total trainable parameter count to accurately estimate fine-tuning memory requirements, and detects the framework version (for example, Keras 2.x versus Keras 3.x) to ensure proper compatibility routing.

For Keras 2.x models incompatible with the primary environment, inspection is delegated to the legacy \texttt{stm32\_legacy} subprocess (see Section~\ref{sec:ai_workflow}).

The extracted architecture is formatted as a human-readable summary and presented to the user, enabling informed decision-making on which layers to freeze, which to replace, and where to introduce regularization.

\subsection{RAG-Improved Modification Planning}

Before soliciting user modifications, the system proactively retrieves best practices tailored to the specific model architecture via the \texttt{retrieve\_best\_practices\_for\_architecture()} function:

\subsubsection{Vectorstore Construction}

The system maintains a ChromaDB vectorstore indexed with documentation from:

\begin{itemize}
  \item Official Keras/TensorFlow guides on transfer learning
  \item Academic papers on efficient fine-tuning techniques (LoRA, adapter layers, etc.)
  \item STMicroelectronics application notes on Edge AI optimization
\end{itemize}

Text is embedded using \texttt{TritonEmbeddings} with the local \texttt{nomic-embed} model and stored locally in \texttt{\textasciitilde/.cache/chroma\_db/architecture\_best\_practices}.

\subsubsection{LLM-Generated Technical Guidance}
\label{sec:llm_generation}

To ensure reliable access to best practices without dependence on external network connectivity or search engine rate limits, the system employs a generative approach using a local LLM (Mistral). Instead of scraping potentially unstable web URLs, the \texttt{\_generate\_and\_cache\_with\_llm()} function prompts an expert system to synthesize a tailored technical guide:

\begin{lstlisting}[language=Python, caption={Local LLM generation of architecture best practices}, label={lst:rag_gen}]
# From workflow5_customization.py

def _generate_and_cache_with_llm(model_name, arch_type):
    llm = ChatTriton(model="mistral", temperature=0.3)
    
    prompt = f"""
    You are an expert embedded AI engineer.
    Provide a concise, bullet-point checklist for fine-tuning 
    {model_name} ({arch_type}) on STM32.
    
    REQUIRED FORMAT:
    * Strategy: [Freeze layers / Retrain behavior]
    * Hyperparams: [LR, Batch Size, Epochs]
    * Quantization: [INT8/float16 recommendations]
    * Constraints: [RAM/Flash usage estimates]
    """
    
    # Generate authoritative guide
    content = llm.invoke(prompt).content
    
    # Cache result in ChromaDB for future instant retrieval
    save_to_vectorstore(content, arch_type)
\end{lstlisting}

This generative approach provides three key advantages. 

First, it ensures deterministic availability by eliminating the ``403 Forbidden'' and timeout errors commonly associated with web scrapers. 

Second, it guarantees structured consistency by enforcing a strict, easily parsable checklist format covering strategy, hyperparameters, and quantization details instead of relying on unpredictable, unstructured HTML. 

Finally, it maximizes search efficiency; while the initial generation may take approximately 15 seconds, all subsequent retrievals are instantaneous thanks to the local vector cache.

\subsubsection{LLM Synthesis}

The top-five most relevant documents, ranked by cosine similarity, are passed to the language model alongside the user's specific architecture description. The LLM then synthesizes a concise recommendation covering critical aspects such as exactly which layers to freeze (e.g., suggesting to freeze the first 80\% of a convolutional backbone to ensure feature extraction stability), the most suitable dropout rates based on the available dataset size, the recommended learning rate ranges for the optimizer, and targeted warnings regarding potential overfitting risks.

This proactive guidance reduces user cognitive load and improves first-attempt success rate for fine-tuning.

Algorithm~\ref{alg:rag_multiquery} formalizes the complete RAG retrieval process with iterative reflection and web search augmentation.

\begin{algorithm}[H]
  \caption{RAG Multi-Query Retrieval with Iterative Reflection}
  \label{alg:rag_multiquery}
  \begin{algorithmic}[1]
    \REQUIRE Model architecture $A$, vectorstore $V$, max loops $L = 3$
    \ENSURE Best practices document $D$
    \STATE $queries \gets \text{LLM\_Generate\_Queries}(A)$ \COMMENT{Generate 6 diverse queries}
    \STATE $ranked\_queries \gets \text{LLM\_Rank\_Queries}(queries, A)$
    \STATE $context \gets \text{""}$
    \FOR{$i \gets 1$ to $|ranked\_queries|$}
    \STATE $weight \gets 1.0 - (i-1) \times 0.05$
    \STATE $docs \gets V.\text{retrieve}(ranked\_queries[i], k=20)$
    \FOR{each $doc \in docs$}
    \STATE $context \gets context + f"\text{Query } i - \text{weight } weight: doc.\text{content}[0:600]"$
    \ENDFOR
    \ENDFOR
    \STATE $summary \gets \text{LLM\_Summarize}(context, A)$
    \STATE $iteration \gets 0$
    \WHILE{$iteration < L$}
    \STATE $gaps \gets \text{LLM\_Reflect}(summary, A)$ \COMMENT{Identify knowledge gaps}
    \IF{$gaps = \text{null}$}
    \RETURN $summary$
    \ENDIF
    \STATE $followup\_query \gets \text{LLM\_Generate\_Followup}(gaps)$
    \STATE $new\_docs \gets V.\text{retrieve}(followup\_query, k=10)$
    \STATE $summary \gets \text{LLM\_Extend}(summary, new\_docs)$
    \STATE $iteration \gets iteration + 1$
    \ENDWHILE
    \RETURN $summary$
  \end{algorithmic}
\end{algorithm}

\subsection{Modification Intent Classification}

The \texttt{ask\_modification\_intent()} node uses a two-stage LLM interaction:

\subsubsection{Stage 1: Binary Decision}

A structured classifier determines whether the user wants to modify the model at all:

\begin{lstlisting}[language=Python, caption={Modification intent detection}, label={lst:intent}]
llm_intent = ChatTriton(model="mistral").with_structured_output(ModificationDecision)

result = llm_intent.invoke([
    SystemMessage(content=modification_decision_instructions),
    HumanMessage(content=user_response)
])

# Output: {wants_modifications: true/false, reasoning: "...", confidence: 0.95}
\end{lstlisting}

If \texttt{wants\_modifications = false}, the workflow bypasses customization and proceeds directly to STEdgeAI code generation.

\subsubsection{Stage 2: Parameter Extraction}

If modifications are desired, a second LLM call extracts structured modification parameters using the \texttt{ModelModification} schema:

\begin{lstlisting}[language=Python, caption={Natural language → structured modification params}, label={lst:mod_extraction}]
# User: "Freeze all layers except the last 5, add 30% dropout, use learning rate 1e-4"

llm_extractor = ChatTriton(model="mistral").with_structured_output(ModelModification)
result = llm_extractor.invoke([
    SystemMessage(content=model_modification_instructions),
    HumanMessage(content=user_response)
])

# Parsed output:
# {
#   "freeze_layers": "all_except_last_5",
#   "dropout_rate": 0.3,
#   "learning_rate": 0.0001,
#   "num_classes": null,
#   "additional_layers": [],
#   "confidence": 0.98
# }
\end{lstlisting}

Supported modifications include:
\begin{itemize}
  \item \texttt{freeze\_layers}: "all", "none", "first\_N", "all\_except\_last\_N", "custom"
  \item \texttt{num\_classes}: Replace output layer for different classification tasks
  \item \texttt{dropout\_rate}: Add dropout regularization (0.0–0.8)
  \item \texttt{learning\_rate}: Override default optimizer LR
  \item \texttt{additional\_layers}: Specify custom layers (e.g., "BatchNormalization", "GlobalAveragePooling2D")
\end{itemize}

\subsection{Standard Fine-Tuning Path}
\label{sec:standard_fine_tuning}

While NNI provides detailed hyperparameter optimization, users often require a faster, single-trial training session for rapid verification. Workflow 5 implements a ``Standard Fine-Tuning'' path (\texttt{fine\_tune\_customized\_model}) that prioritizes rapid feedback and resource-aware stability.

\subsubsection{Anti-Overfitting Logic}

Large pre-trained models (e.g., MobileNetV2) are highly prone to catastrophic forgetting and severe overfitting when applied to small or specialized datasets like Fruit-360. To effectively mitigate this, the system applies three automated strategies. 

First, it employs a dynamic sampling escalation strategy that caps dataset usage (e.g., at 5000 images) to balance class representation while preventing memory exhaustion and excessive latency. 

Second, the system proactively inserts a regularization \texttt{Dropout} layer (with a default rate of 0.3) during the final layer replacement phase; importantly, if the user explicitly specifies a custom dropout rate in their natural language prompt, the system overrides the safety default in favor of the user's instruction. Lastly, the system implements epoch capping, strictly limiting standard training runs to 10 epochs to act as a rapid ``sanity check'' before triggering any computationally expensive optimization loops.

\subsubsection{Real-time Observability and Fault Tolerance}

To effectively connect between headless background subprocesses and the interactive agent web UI, the fine-tuning execution engine provides high-fidelity, real-time observability. It utilizes asynchronous log streaming via \texttt{subprocess.Popen} with non-blocking pipe selection to instantly stream training metrics, such as epochs, loss, and accuracy, directly to the user's console. To maintain stability, the engine features adaptive timeouts, intelligently adjusting the maximum allowed execution window (e.g., 1200s) based on the model's complexity and the estimated dataset size. 

Also, the system includes a class mismatch rescue mechanism; if it detects that the loaded dataset contains a different number of classes than the pre-trained model (for instance, 132 versus 1000), it automatically performs an architectural adjustment by replacing the final \texttt{Dense} layer and dynamically resizing the output tensor before commencing any training.

\subsection{NNI Code Generation: Architecture and Prompt Engineering}

The core technical contribution of this workflow is the \textit{autonomous generation of complete NNI experiment code} from high-level specifications. Unlike rigid, template-based code generators that rely on simple fill-in-the-blank patterns, this system employs a large language model (GPT-OSS:20b with an 8192-token context window) to dynamically synthesize production-ready Python scripts. These scripts are solid enough to handle complex edge cases such as sparse versus categorical label formats, input shape mismatches requiring dynamic image resizing, and class count discrepancies between the base model and the target dataset. 

Also, the generated code gracefully manages memory constraints on consumer hardware through dataset subsetting and memory-mapping, resolves port conflicts for the NNI Web UI, and implements reliable fallbacks for missing NNI API methods across different library versions.

\subsubsection{Prompt Engineering Strategy}

The \texttt{generate\_nni\_experiment()} function in \texttt{nni\_optimization/generator.py} constructs a highly detailed, 480-line prompt designed to ensure accurate generation. The prompt heavily features context injection, providing the LLM with exact model metadata (such as name, tensor shapes, layer count, and file path) along with detailed dataset specifications. It enforces explicit constraints, strictly prohibiting the use of placeholder paths, the omission of required files, and the hallucination of non-existent API methods. To guide the output structure, the prompt includes detailed working code templates for both \texttt{manager.py} and \texttt{trial.py}, annotated with inline comments that explain computationally critical sections. 

Also, it embeds specific error prevention rules regarding known NNI API inconsistencies (for example, explicitly warning that \texttt{experiment.get\_best\_trial()} does not exist in certain versions). 

Finally, it mandates a strict output format specification, requiring the LLM to use explicit \texttt{\# FILE:} markers to successfully generate multi-file outputs.

A representative excerpt from the prompt:

\begin{verbatim}
CRITICAL REQUIREMENTS:
- YOU MUST GENERATE EXACTLY 2 FILES: manager.py AND trial.py
- Use the exact format shown below with # FILE: markers
- DO NOT use keras.datasets.load_data(). ALWAYS load data from {dataset_path}.
- IMPORTANT: If data labels (y_train) have shape (N, 1), they are SPARSE. 
   Convert them to categorical (one-hot) before training.

FILE 1: manager.py
- Configure NNI experiment (TPE tuner, maximize mode)
- Define search space (learning_rate, batch_size, optimizer, freeze_backbone)
- Set trial_command = "python trial.py"
- Launch experiment with experiment.run(port=8080, wait_completion=True)
- AFTER experiment: Find the best trial USING ONLY experiment.list_trial_jobs()
- Trigger RETRAIN: Run "trial.py" as subprocess with env RETRAIN_MODE='true'

FILE 2: trial.py
- Get parameters (from NNI_PARAMS if RETRAIN_MODE, else nni.get_next_parameter())
- Load model and data from provided paths (USE EXACT PATHS: {model_path} and {dataset_path})
- MEMORY OPTIMIZATION: Load data with mmap_mode='r' to avoid RAM overload
- DATASET SUBSET: Use only 50% of the dataset to save RAM
- Adapt model final layer if classes mismatch
- Train model for 3 epochs
- Report accuracy to NNI (if not retrain) or Save 'best_model.h5' (if retrain)

Critical Rules:
- DO NOT USE experiment.export_top_trial() OR experiment.get_best_trial() - THEY DO NOT EXIST.
- ALWAYS USE experiment.list_trial_jobs() to find the best trial manually.
- Use exact paths from context: {dataset_path} and {model_path}
\end{verbatim}

The system temperature is set to 0 for deterministic code generation.

\subsubsection{Multi-File Parser with Regex Markers}

The LLM outputs a single concatenated string containing both files, delimited by \texttt{\# FILE: filename.py} markers. The \texttt{extract\_files\_from\_code()} function parses this using a permissive regex pattern:

\begin{lstlisting}[language=Python, caption={Multi-file code extraction}, label={lst:file_extraction}]
file_pattern = re.compile(r'^\#\s*FILE:\s*([a-zA-Z0-9_\-\.]+)', re.MULTILINE | re.IGNORECASE)
file_matches = list(file_pattern.finditer(code))

files = {}
for i, match in enumerate(file_matches):
    filename = match.group(1)
    start = match.end()
    
    # Find next file marker or end of string
    if i + 1 < len(file_matches):
        end = file_matches[i + 1].start()
    else:
        end = len(code)
    
    file_content = code[start:end].strip()
    files[filename] = file_content

# Result: {"manager.py": "import nni\n...", "trial.py": "import sys\n..."}
\end{lstlisting}

The parser is intentionally designed to robustly handle common minor variations in the LLM's raw output string. It effectively ignores optional spacing around the hash or \texttt{FILE:} markers, securely strips away any markdown code fences (like \texttt{```python}) that might wrap the file content, and correctly identifies and discards trailing comments mistakenly appended after the filename.

If no markers are found, the entire output is saved as \texttt{main.py} (fallback for single-file tasks).

\subsubsection{Generated File Structure}

The typical output structure for an NNI experiment:

\paragraph{\texttt{manager.py}}
This file orchestrates the entire hyperparameter search:

\begin{itemize}
  \item \textbf{Environment Setup}: Fixes CUDA library paths for Conda environments (lines 192-196)
  \item \textbf{Search Space Definition}: Defines hyperparameter ranges (learning\_rate, batch\_size, optimizer, freeze\_backbone)
  \item \textbf{Experiment Configuration}: Sets tuner (TPE/Random), max trials, trial concurrency, GPU usage
  \item \textbf{Port Hunting}: Scans for free ports (8080–8100) to avoid conflicts (lines 226-235)
  \item \textbf{Experiment Execution}: Launches trials with \texttt{experiment.run(wait\_completion=True)}
  \item \textbf{Best Trial Extraction}: Manually iterates \texttt{experiment.list\_trial\_jobs()} to find highest-accuracy trial (lines 244-253)
  \item \textbf{Auto-Retrain}: Re-runs \texttt{trial.py} with best hyperparameters in \texttt{RETRAIN\_MODE} to save final model
\end{itemize}

\paragraph{\texttt{trial.py}}
This file implements a single training trial:

\begin{itemize}
  \item \textbf{Mode Detection}: Checks \texttt{RETRAIN\_MODE} environment variable to determine behavior
  \item \textbf{Parameter Loading}: Retrieves hyperparameters from NNI API or environment variable
  \item \textbf{Data Loading}: Memory-mapped NumPy array loading (\texttt{mmap\_mode='r'}) to avoid RAM exhaustion
  \item \textbf{Dataset Subsetting}: Uses only 50\% of data to fit in consumer RAM (lines 373-381)
  \item \textbf{Label Format Fix}: Detects sparse labels (shape \texttt{(N, 1)}) and converts to categorical (one-hot)
  \item \textbf{Adaptive Output Layer}: Replaces final Dense layer if class count mismatches (lines 400-408)
  \item \textbf{Transfer Learning Logic}: Freezes backbone layers if \texttt{freeze\_backbone = True}
  \item \textbf{Dynamic Preprocessing}: Resizes input images to match model's expected shape via TensorFlow Dataset API
  \item \textbf{Data Augmentation}: Applies random flips, rotations, zoom during training
  \item \textbf{Result Reporting}: Sends \texttt{val\_accuracy} to NNI or saves \texttt{best\_model.h5} (retrain mode)
\end{itemize}

\subsubsection{Memory Management and GPU Offloading}

To prevent GPU OOM errors when transitioning from LLM inference to NNI training, the system implements explicit resource cleanup:

\begin{lstlisting}[language=Python, caption={Triton model unloading before NNI}, label={lst:gpu_cleanup}]
def force_unload_triton(model_name: str = "mistral"):
    """
    Force Triton to unload the model from GPU to free up VRAM for TensorFlow/NNI.
    Sends a repository management request.
    """
    url = f"http://localhost:8000/v2/repository/models/{model_name}/unload"
    
    req = urllib.request.Request(url, method='POST')
    try:
        with urllib.request.urlopen(req) as response:
            logger.info(f"Triton model '{model_name}' unloaded to free VRAM for NNI.")
    except Exception as e:
        logger.error(f"Failed to unload Triton model: {e}")
\end{lstlisting}

This function is called in the \texttt{finally} block of \texttt{generate\_nni\_experiment()} to guarantee VRAM release even if code generation fails.

\subsection{Hyperparameter Optimization Execution}

Once the NNI scripts are generated and saved to \texttt{./nni\_experiment/}, the \texttt{optimize\_hyperparameters\_with\_nni()} function executes them:

\begin{enumerate}
  \item \textbf{Validation}: Checks that both \texttt{manager.py} and \texttt{trial.py} exist and are syntactically valid Python
  \item \textbf{Subprocess Launch}: Spawns \texttt{python manager.py} in a separate process to avoid blocking the LangGraph event loop
  \item \textbf{Output Streaming}: Captures stdout/stderr in real-time for user feedback
  \item \textbf{Completion Detection}: Waits for subprocess termination (blocking until experiment finishes)
  \item \textbf{Model Retrieval}: Locates \texttt{best\_model.h5} in the NNI output directory and stores path in state
\end{enumerate}

The subprocess execution isolates NNI's environment from the LangGraph runtime, preventing library conflicts and allowing independent GPU scheduling.

Algorithm~\ref{alg:nni_retrain} formalizes the complete NNI workflow including automatic best-trial extraction and retraining.

\begin{algorithm}[H]
  \caption{NNI Experiment with Automatic Best-Trial Retraining}
  \label{alg:nni_retrain}
  \begin{algorithmic}[1]
    \REQUIRE Model path $M$, dataset path $D$, search space $S$, max trials $N$
    \ENSURE Retrained best model saved to disk
    \STATE $experiment \gets \text{NNI\_Create}(\text{tuner}=\text{"TPE"}, S, N)$
    \STATE $experiment.\text{run}(\text{wait\_completion}=\text{true})$
    \STATE \COMMENT{Extract best trial manually (stable API-agnostic method)}
    \STATE $trials \gets experiment.\text{list\_trial\_jobs}()$
    \STATE $valid\_trials \gets \{t \in trials \mid t.\text{status} = \text{"SUCCEEDED"}\}$
    \IF{$|valid\_trials| = 0$}
    \STATE \textbf{error} "No successful trials"
    \ENDIF
    \STATE $best\_trial \gets \arg\max_{t \in valid\_trials} t.\text{finalMetricData}[0].\text{data}$
    \STATE $best\_params \gets best\_trial.\text{hyperParameters}[-1].\text{parameters}$
    \STATE \COMMENT{Retrain with best hyperparameters}
    \STATE $env \gets \{\text{"RETRAIN\_MODE"}: \text{"true"}, \text{"NNI\_PARAMS"}: \text{JSON}(best\_params)\}$
    \STATE $\text{subprocess.run}([\text{"python", "trial.py"}], \text{env}=env)$
    \STATE $best\_model \gets \text{Load}(\text{"best\_model.h5"})$
    \RETURN $best\_model$
  \end{algorithmic}
\end{algorithm}

\subsection{Validation and Model Persistence}

After NNI completes, the \texttt{validate\_customized\_model()} node loads \texttt{best\_model.h5} and evaluates it on a held-out test set. If validation accuracy meets the user's threshold (configurable, default 0.7), the model is persisted with metadata:

\begin{lstlisting}[language=Python, caption={Model persistence with metadata tracking}, label={lst:model_save}]
metadata = {
    "base_model": state.model_path,
    "customization_date": datetime.now().isoformat(),
    "nni_search_space": state.nni_search_space,
    "best_hyperparameters": state.best_hyperparameters,
    "validation_accuracy": final_accuracy,
    "framework_version": "LangGraph-STM32-MLOps v1.0"
}

model.save(output_path)
with open(output_path.replace('.h5', '_metadata.json'), 'w') as f:
    json.dump(metadata, f, indent=2)
\end{lstlisting}

This metadata enables reproducibility and audit trails for deployed models.
